\subsection{Estado de conflicto y situación inicial}
Una vez identificados los actores y el usuario principal, se vuelve necesario representar cómo opera actualmente el sistema y dónde se produce el conflicto que da origen al proyecto. La diagramación de la situación inicial permite describir el estado actual del problema, visibilizando flujos, restricciones, puntos de fricción y efectos sobre los distintos involucrados. De esta manera, no solo se enuncia una problemática, sino que se estructura de forma clara el escenario desde el cual emerge la necesidad de intervención.\\

\subsubsection*{Situación inicial} ~

En el contexto actual, la información clínica del paciente es registrada por el equipo de salud en la ficha clínica electrónica oficial del hospital. Sin embargo, ese registro está diseñado para almacenar datos transaccionales y legales, no para capturar el contexto cotidiano del caso. Los eventos menores, los acuerdos de cuidado, las observaciones de turno o los detalles de seguimiento que no justifican un registro formal son precisamente los que resultan más críticos para la continuidad asistencial, y son también los que con mayor frecuencia se pierden.

El modelado de la situación actual o AS-IS se realiza para representar de manera ordenada el funcionamiento presente del sistema antes de proponer cualquier solución. Esta herramienta permite identificar cómo circula actualmente la información, qué actores intervienen en ese flujo, qué barreras aparecen en el proceso y dónde se generan quiebres entre la información disponible y la información efectivamente traspasada al equipo entrante. Su utilidad radica en evidenciar que el problema no es la inexistencia de información clínica, sino la ausencia de una capa que la capture en el momento correcto, la organice de forma comprensible y la ponga disponible para el equipo que la necesita.

Paciente hospitalizado $\rightarrow$ personal clínico genera eventos asistenciales $\rightarrow$ parte se registra en ficha oficial, parte en canales informales, parte no se registra $\rightarrow$ cambio de turno $\rightarrow$ médico saliente realiza traspaso verbal parcial $\rightarrow$ médico entrante toma decisiones con contexto incompleto $\rightarrow$ aumenta riesgo clínico y tiempo de contextualización $\rightarrow$ se incrementa la presión sobre el equipo asistencial.\\

\subsubsection*{Estado del conflicto} ~

La definición del estado de conflicto tiene como objetivo explicitar la tensión central que estructura la problemática del proyecto, identificando las fuerzas contrapuestas que mantienen el problema sin resolución en el sistema actual.

El conflicto principal surge en la intersección de tres tensiones. En primer lugar, el personal asistencial necesita registrar el contexto clínico del paciente de forma rápida y sin fricción, pero los sistemas disponibles están diseñados para registros formales y transaccionales que requieren tiempo y protocolo. En segundo lugar, el equipo médico entrante necesita acceder al contexto actualizado del caso al momento del cambio de turno, pero en los sistemas descritos ese contexto no aparece de forma estructurada ni accesible. En tercer lugar, la institución requiere trazabilidad de los eventos asistenciales para cumplir con sus obligaciones normativas y de calidad, pero los canales informales que el personal utiliza para compensar las limitaciones del sistema oficial no dejan registros auditables.

La problemática no es solo informacional, sino también operativa, cultural y de diseño de sistemas. Existe información, pero no existe una capa que la capture en el momento correcto, la organice de forma comprensible y la ponga disponible para el equipo entrante con el objetivo de reducir la duplicación de trabajo. La oportunidad de diseño se encuentra en esa brecha: desarrollar una capa de trazabilidad complementaria que registre el contexto cotidiano del paciente, lo organice cronológicamente y lo haga accesible al equipo clínico en el momento en que más lo necesita, con el objetivo de no reemplazar la ficha oficial ni aumentar la carga operativa del personal.
